\documentclass[a4paper,draft]{article}

\usepackage{csquotes}
\usepackage{hyperref}
\usepackage[british]{babel}
\usepackage{booktabs}
\usepackage{subfiles}
\usepackage{amsmath}
\usepackage{cleveref}
\usepackage{fixme}

%% References
\usepackage[
    style=numeric-comp,
    backend=biber
]{biblatex}
\addbibresource{../library.bib}

%% Acronyms
\usepackage[automake]{glossaries-extra}
\setabbreviationstyle[acronym]{long-short}
\makeglossaries

% Apparently the only way to prevent all the acronyms from generating spelling errors is to 
% disable ltex on the entire file.
% LTeX: enabled=false

\newabbreviation{ai}{AI}{Artificial Intelligence}
\newabbreviation{ar}{AR}{Augmented Reality}
\newabbreviation{asf}{ASF}{Amplitude Stereo Features}
\newabbreviation{cac}{CAC}{Cut And Clicks}
% \newabbreviation{cd}{CD}{Cochleagram Difference}
\newabbreviation{cd}{CD}{Compact Disc}
\newabbreviation{dft}{DFT}{Discrete Fourier Transform}
\newabbreviation{dirac}{DirAC}{Directional Audio Coding}
\newabbreviation{dl}{DL}{Deep Learning}
\newabbreviation[longplural=Directions of Arrival]{doa}{DOA}{Direction of Arrival}
\newabbreviation{dsp}{DSP}{Digital Signal Processing}
\newabbreviation{erb}{ERB}{Equivalent Rectangular Bandwidth}
\newabbreviation{fa}{FA}{Frequency-Azimuth}
\newabbreviation{faa}{FAA}{Frequency-Azimuth Analysis}
\newabbreviation{fft}{FFT}{Fast Fourier Transform}
\newabbreviation{fir}{FIR}{finite impulse response}
\newabbreviation{fmttm}{FMTTM}{Fly Me To The Moon}
\newabbreviation{genai}{Gen-AI}{Generative Artificial Intelligence}
\newabbreviation{iamf}{IAMF}{Immersive Audio Model and Formats}
\newabbreviation{daw}{DAW}{Digital Audio Workstation}
\newabbreviation{dj}{DJ}{Disc Jockey}

\newabbreviation{rq}{RQ}{research question}

\newabbreviation{iid}{IID}{Interaural Intensity Difference}
\newabbreviation{ipd}{IPD}{Interaural Phase Difference}
\newabbreviation{icpd}{ICPD}{Inter-channel Phase Difference}
\newabbreviation{ccpd}{CCPD}{Cross-channel Phase Difference}
\newabbreviation{ild}{ILD}{Interaural Level Difference}
\newabbreviation{ilr}{ILR}{Interaural Level Ratio}
\newabbreviation{cld}{CLD}{Channel Level Difference}
\newabbreviation{ccad}{CCAD}{Cross-channel Amplitude Difference}
\newabbreviation{icld}{ICLD}{Inter-channel Level Difference}
\newabbreviation{itd}{ITD}{Interaural Time Difference}
\newabbreviation{ictd}{ICTD}{Inter-channel Time Difference}
\newabbreviation{iacc}{IACC}{Interaural Cross Correlation}
\newabbreviation{iccc}{ICCC}{Inter-channel Cross Correlation}
\newabbreviation{icc}{ICC}{Inter-channel Correlation}
\newabbreviation{ic}{IC}{Inter-channel Coherence}
\newabbreviation{iac}{IAC}{Interaural Coherence}
\newabbreviation{itf}{ITF}{Inter-channel Transfer Function}
\newabbreviation{lev}{LEV}{Listener Envelopment}
\newabbreviation{dsd}{DSD}{Directional Loudness Distortion}
\newabbreviation{bms}{BMS}{Bivariate Mixture Space}
\newabbreviation{spl}{SPL}{Sound Pressure Level}
\newabbreviation{iir}{IIR}{Infinite Impulse Response}


\newabbreviation{ishst}{ISHST\textsuperscript{**}}{I Saw Her Standing There\textsuperscript{**}}
\newabbreviation{kks}{KKS}{The Knowledge Foundation}
\newabbreviation{lb}{LB}{Laser Beam}
\newabbreviation{lfe}{LFE}{Low Frequency Effects}
\newabbreviation{llm}{LLM}{Large Language Model}
\newabbreviation{lufs}{LUFS}{Loudness Units Full Scale}
\newabbreviation{mirex}{MIREX}{Music Information Retrieval Evaluation eXchange}
\newabbreviation{mir}{MIR}{Music Information Retrieval}
\newabbreviation{ml}{ML}{Machine Learning}
\newabbreviation{mtg}{MTG}{Music Technology Group}
\newabbreviation{m}{M}{Manifesto}
\newabbreviation{pc}{PC}{Panning Coefficients}
\newabbreviation{ph}{PH}{Panning Histogram}
\newabbreviation{pi}{PI}{Panning Index}
\newabbreviation{picoc}{PICOC}{Population, Intervention, Comparison, Outcome, Context}
\newabbreviation{psf}{PSF}{Phase Stereo Features}
\newabbreviation{qm}{QM}{Qualitative Model}
\newabbreviation{rms}{RMS}{Root Mean Squared}
\newabbreviation{rm}{RM\textsuperscript{*}}{'Round Midnight\textsuperscript{*}} 
\newabbreviation{simpl}{SIMPL}{stereo-input mixed-to-panned level}
\newabbreviation{sip}{SIP}{Stereo Image Profile}
\newabbreviation{slr}{SLR}{Systematic Literature Review}
\newabbreviation{spf}{SPF}{Stereo Panning Features}
\newabbreviation{sps}{SPS}{Stereo Panning Spectrum}
\newabbreviation{ssps}{SSPS}{Spectral Stereo Phase Spread}
\newabbreviation{stft}{STFT}{Short-time Fourier Transform}
\newabbreviation{tf}{TF}{time-frequency}
\newabbreviation{titt}{TITT}{This Is The Thing}
\newabbreviation{vr}{VR}{Virtual Reality}
\newabbreviation{peaq}{PEAQ}{Objective Measurements of Perceived Audio Quality}
\newabbreviation{peaqmc}{PEAQ-MC}{Objective Measurements of Perceived Audio Quality - Multi Channel}
\newabbreviation{cnn}{CNN}{Convolutional Neural Network}

\newabbreviation{ica}{ICA}{Independent Component Analysis}
\newabbreviation{pca}{PCA}{Principal Component Analysis}
\newabbreviation{pqmf}{PQMF}{Pseudo Quadrature Mirror Filter}
\newabbreviation{mpeg}{MPEG}{Moving Picture Experts Group}
\newabbreviation{3gpp}{3GPP}{3rd Generation Partnership Project}

\newabbreviation{asw}{ASW}{Apparent Source Width}
\newabbreviation{fwhm}{FWHM}{Full Width Half Maximum}
\newabbreviation{lef}{LEF}{Lateral Energy Fraction}
\newabbreviation{iqr}{IQR}{Interquartile Range}
\newabbreviation{amt}{AMT}{Auditory Modelling Toolbox}
\newabbreviation{qestral}{QESTRAL}{Quality Evaluation of Spatial Transmission and Reproduction using an Artificial Listener}
\newabbreviation{bcc}{BCC}{Binaural Cue Coding}
\newabbreviation{ps}{PS}{Parametric Stereo}
\newabbreviation{saoc}{SAOC}{Spatial Audio Object Coding}
\newabbreviation{gccphat}{GCC-PHAT}{Generalized Cross Correlation with Phase Transform}
\newabbreviation{mp3}{MP3}{MPEG layer 3}
\newabbreviation{aac}{AAC}{Advanced Audio Coding}
\newabbreviation{ivas}{IVAS}{Immersive Voice and Audio Services}
\newabbreviation{sofa}{SOFA}{Spatially Oriented Format for Acoustics}

\newabbreviation{ir}{IR}{Impulse Response}
\newabbreviation{hrir}{HRIR}{Head-Related Impulse Response}
\newabbreviation{hrtf}{HRTF}{Head-Related Transfer Function}
\newabbreviation{midi}{MIDI}{Musical Instrument Digital Interface}
\newabbreviation{vi}{VI}{Virtual Instrument}
\newabbreviation{mushra}{MUSHRA}{MUlti Stimulus Hidden Reference and Anchor}
\newabbreviation{bsseval}{BSS\_EVAL}{Blind Source Separation Evaluation}
\newabbreviation{psd}{PSS}{Perceptual Spectral Difference}
\newabbreviation{bamq}{BAM-Q}{Binaural-Auditory-Model-based Quality}
\newabbreviation{tff}{TFF}{Time-Frequency Filter}
\newabbreviation{pco}{PCO}{Population, Context, Outcome}
\newabbreviation{wos}{WOS}{Web Of Science}
\newabbreviation{ss}{SS}{Semantic Scholar}
\newabbreviation{cpu}{CPU}{Central Processing Unit}
\newabbreviation{flop}{FLOP}{Floating Point Operation}

\newabbreviation{csv}{CSV}{Comma Separated Variable}
\newabbreviation{tsv}{TSV}{Tab Separated Variable}
\newabbreviation{api}{API}{Application Programming Interface}
\newabbreviation{doi}{DOI}{Digital Object Identifier}
\newabbreviation{url}{URL}{Universal Resource Locator}
\newabbreviation{prisma}{PRISMA}{Preferred Reporting Items for Systematic reviews and Meta-Analyses}
\newabbreviation{adress}{ADRess}{Azimuth Discrimination and Resynthesis}
\newabbreviation{agr}{AGR}{Automatic Genre Recognition}
\newabbreviation{som}{SOM}{Self Organising Map}
\newabbreviation{iresap}{IRESAP}{Information Retrieval in Embedded Systems for Audiovisual Artistic Processes}
\newabbreviation{mdu}{MDU}{Mälardalen Universitet}
\newabbreviation{kmh}{KMH}{Kungliga Musikhögskolan}




% datasets
\newabbreviation{boma}{BOMA}{{Bang \& Olufsen Music for Archimedes}}
\newabbreviation{uiowamis}{UIowaMIS}{University of Iowa Musical Instrument Samples}
% \newabbreviation{}{}{Bach10}
\newabbreviation{btpab}{BTPAB}{ByTeDance paired ambisonic and binaural}
\newabbreviation{dsd100}{DSD100}{Demixing Secret Database}
\newabbreviation{asascbrm}{ASASCBRM}{Database for Automatic Spatial Audio Scene Classification in Binaural Recordings of Music}
\newabbreviation{daoc}{DAOC}{Denon anechoic orchestral collection}
% \newabbreviation{}{}{FAIR-Play}
\newabbreviation{fma}{FMA}{Free Music Archive}
% \newabbreviation{}{}{MedleyDB}
% \newabbreviation{}{}{MedleyDB 2.0}
% \newabbreviation{}{}{MUSDB18}
% \newabbreviation{}{}{MUSDB18-HQ}
\newabbreviation{quasi}{QUASI}{QUaero Audio SIgnals}
\newabbreviation{rwc}{RWC}{Real World Computing}
\newabbreviation{sqam}{SQAM}{EBU Sound Quality Assessment Material}
% \newabbreviation{}{}{Slakh}
\newabbreviation{urmp}{URMP}{University of Rochester Multi-Modal Music Performance}
% \newabbreviation{}{}{ccMixter}
\newabbreviation{arsso}{ARSSO}{Anechoic Recording System for Symphony Orchestra}
\newabbreviation{rioossr}{RIOOSSR}{Recordings of Italian opera orchestra and soloists in a silent room}
\newabbreviation{msss}{MSSS}{Mixing Secrets for The Small Studio}
\newabbreviation{tss}{TSS}{Telefunken Studio Sessions}
\newabbreviation{mtgj}{MTGJ}{Music Technology Group Jamendo}
% \newabbreviation{}{}{Audio and Acoustic Signal Processing Technical Committee Online Resources}
\newabbreviation{maestro}{MAESTRO}{MIDI and Audio Edited for Synchronous TRacks and Organization}
% \newabbreviation{}{}{FiloSax}
% \newabbreviation{}{}{Gamelan Music Dataset}
% \newabbreviation{}{}{mobygratis}


% Journals
\newabbreviation{jaes}{JAES}{Journal of the Audio Engineering Society}
\newabbreviation{eurasip}{EURASIP}{European Association for Signal Processing}

\title{Licentiate Proposal - Analysis of stereo image in vast repositories of recorded music}

\author{Richard Mitic}

\begin{document}

\maketitle


\begin{abstract}
\subfile{abstract}
\end{abstract}

\pagebreak


\listoffixmes

\section{Papers included in the thesis}
\label{sec:papers}

\begin{refsection}

\subsection{Paper A}
\label{paper:A}
\fullcite{franklinDevelopmentWebBasedRealTime2025}
\\\\
\textbf{Status} -- published March 2025

\subsection{Paper B}
\label{paper:B}

\fullcite{miticObjectiveComparisonPanning2025}
\\\\
\textbf{Status} -- published September 2025

\subsection{Paper C}
\label{paper:C}
\fullcite{miticPerceptualEvaluationBinaurally2026}
\\\\
\textbf{Status}
\begin{itemize}
    \item Submitted to \href{https://dl.acm.org/journal/tomm}{ACM Transactions on Multimedia Computing Communications and Applications} October 2025 - \textbf{Rejected - not good fit for journal}.
    \item Submitted to \href{https://acta-acustica.edpsciences.org/}{Acta Acustica} February 2025
\end{itemize}
% \\
\textbf{Plan} -- Await peer review.

\subsection{Paper D}
\label{paper:D}
\fullcite{miticSystematicLiteratureReview2026}
\\\\
\textbf{Status}
\begin{itemize}
\item Submitted to \href{https://asmp-eurasipjournals.springeropen.com/}{EURASIP Journal on Audio, Speech, and Music Processing} on 2025-12-16.
\item Feedback received 2026-02-02. Changes required -- see below.
\end{itemize}
\textbf{Plan} -- Work through required changes. Deadline is 2026-04-03.
\begin{itemize}
    \item Update search results for the period between original search and submission.
    \item Get a colleague to double-check a subset of the original screening.
    \item Add more, smaller diagrams, and break up existing big diagrams.
    \item Reorganize text to make relationship between discovered papers more obvious.
    \item Minor clarifications and typos.
\end{itemize}


\end{refsection}

\pagebreak

\section{Introduction}

It is a widely accepted truth in the world of music production and recording that spatial information in audio is crucial to  listeners' sense of immersion and connection with the music. To capture that spatiality, stereophonic recording and reproduction has been the standard format since it gained popularity in the 1950s \cite{lanierHIFITWOCHANNELCOMMOTION1957}. The spatial analysis of subsequent decades worth of stereo recordings presents a unique opportunity to study human musicality, as each recording reflects and encodes the creative choices made by the musicians and recordists involved.

Within the music community, it is common to speak about spatiality in a qualitative sense, and there have been several attempts to translate these discussions in a quantitive domain \cite{tzanetakisStereoPanningFeatures2007,dewhirstModellingPerceivedSpatial2008,sarroffMeasurementsSpaciousnessStereophonic2008,tardieuProductionEffectAudio2011}. All of these studies were performed in an era when musical datasets were hand-curated/-labelled, fairly small (1,000 - 2,000 tracks), and -- most notably -- private. However, in recent years much larger public musical corpora have become available, e.g. MTG-Jamendo \cite{bogdanovMTGJamendoDatasetAutomatic2019} and \gls{fma} \cite{defferrardFMADatasetMusic2017}, which contain approximately 55,000 and 100,000 tracks respectively. In addition, computing techniques in \gls{ml} and \gls{dl} have progressed to the point where it is now possible to learn extremely high-level features on these vast datasets, thereby making sense of the human creativity represented in the recordings. 

A typical use case for spatial descriptors that has emerged both in academia and industry is classification, and in particular, auto-tagging of musical genre/performer \cite{tzanetakisStereoPanningInformation2010,ziemerGoniometersArePowerful2023}. Such features can be used in music similarity and recommendation tasks \cite{gomezMusicSimilaritySystems2008,ziemerNovelRecordingStudio2020}. It has also been shown that stereo field classification of subsections within a larger work can act as a precursor to other tasks such as audio enhancement, upmixing, and source separation \cite{harmaStereoAudioClassification2011,harmaClassificationTimefrequencyRegions2011}. The latter is very relevant to the audio distribution industry, which is currently experimenting with newer audio formats such as ambisonic and object-based audio. While these new formats offer exciting opportunities for both distributors and listeners, there remains a huge back-catalogue of stereophonic recordings that could benefit from a high-level understanding of their spatial content. 

Historically, binaural audio (which is a subclass of stereo) has been of particular interest because its relevance to human hearing. Thus, it is extremely widely studied and the underlying fundamental techniques are well-established. There remains however, a gap in current research regarding direct analysis of raw music recordings, and in particular, describing the stereo image of said recordings at a high level. This thesis aims to fill that gap. It will firstly demonstrate use cases where the spatial analysis of stereo audio has the opportunity to solve real-world problems in audio transmission, dissemination, and reproduction. Accordingly, it will gather existing methods to describe stereo image and assess their abilities to detect the spatial aspects of a variety of musical audio sources.

\section{Research Questions}
\label{sec:request_questions}

The overarching research question can be formed as:

\begin{quote}
How do music producers use stereo field?
\end{quote}

\noindent This can be broken into several smaller problems:

\begin{enumerate}
\item How does the stereo image of a recording relate to real-world musical knowledge?
\item What methods can be used to describe the stereo image of music recording?
\end{enumerate}

\fxfatal{Expand notes of how each paper related to the research questions}
Auxtrument paper - real use case where stereo field and reproduction thereof is important. One conclusion was that network bandwidth and latency was a major factor. Audio quality needed to remain high so that spatial qualities were preserved, but that resulted in high bandwidth. It would've been great to have a standard audio format that focused on keeping spatial information intact.

BTH paper - Another real use case where stereo image profile can be used as a precursor to a different task. Speech was panned binaurally and mixed with music. Participants were asked to rate the pleasantness of each variant. The results showed that the default centre position was not always preferred, although more data would need to be collected to directly link the SIP with the best panning position. 

SLR - gather all existing methods to generate a SIP. It breaks each one down into its fundamental components and draws links between them, finally culminating in a framework for creating new SIPs. It covers both oldskool DSP methods and newer \gls{dl} methods.

DAFx paper - Compare several very relevant SIP methods on real music to see what they can detect. It assigns a score to each one denoting how well they are able to separate some tracks that were hand-picked to represent different used of stereo field.

\section{State of the art}

\Citeauthor{bergSystematicEvaluationPerceived2003} \cite{bergSystematicEvaluationPerceived2003} specifie a long list of spatial attributes perceived by listeners, but three major ones often referred to in the literature are \enquote{ensemble source width} (sometimes called \enquote{apparent source width}), \enquote{extent of immersion} (sometimes called \enquote{listener envelopment}), and \enquote{extent of reverberation} \cite{soulodreNewObjectiveMeasures2004,sarroffMeasurementsSpaciousnessStereophonic2008,sarroffComputationalModelPerceived2011}. \cite{ziemerUsingPsychoacousticModels2016} and \cite{ziemerSourceWidthMusic2017a} define a scalar measurement of \enquote{spaciousness} (which is very closely related to source width). Other lower-level scalar features exist, such as the \gls{simpl} feature which is used in \cite{chenMixedStereoAudio2014} to aid the classification of speech vs music. 

The above perceptual qualities are often inferred using objective measurements, which aim to describe the physical properties of an audio signal (usually in higher dimensionality) than the perceptual scalars.
\cite{avendanoFrequencyDomainTechniques2002,avendanoFrequencydomainSourceIdentification2003} are seminal works describing the \enquote{panning index} -- a description of amplitude panning for each \gls{tf} bin in a stereo spectrogram. \cite{avendanoFrequencyDomainTechniques2002} also presents the \enquote{ambience index} -- a measure of inter-channel correlation in \gls{tf} bins. While they target re-synthesis and upmix as use cases, the underlying techniques have become the basis for many stereo image descriptors since. 
\cite{barrySoundSourceSeparation2004} is another seminal work presenting the \gls{adress} algorithm for separating individual musical sources from a mix. In doing so it describes a method for determining the panning position of individual \gls{tf} bins which has been used since to create spatial descriptors. 
\cite{kimQuantifyingAlbumEffect2006} states that analysis of a recording's stereo field can be useful for classification, but does not give any techniques to measure it.
\cite{tzanetakisStereoPanningFeatures2007,tzanetakisStereoPanningInformation2010} reuse the panning index from \cite{avendanoFrequencyDomainTechniques2002} (renamed to \enquote{\gls{sps}}) and derive a set of \gls{spf} that are used for genre classification and \gls{mir}.
Similarly, \cite{sarroffMeasurementsSpaciousnessStereophonic2008} details a method to produce stereo image descriptors centred around a histogram of spatial energy, but uses the frequency-azimuth plane from \cite{barrySoundSourceSeparation2004} as its basis.
Another histogram method, this time based on the \gls{sps} from \cite{tzanetakisStereoPanningFeatures2007} is described in \cite{gomezMusicSimilaritySystems2008,guausAudioContentProcessing2009}, citing measurement of music similarity as the primary use case.
In a similar vein, and citing classification as a use case, \cite{tardieuProductionEffectAudio2011} combines cochleagram differences with an analysis of phase differences between channels to form two feature sets - the \enquote{\gls{asf}} and \enquote{\gls{psf}}.

\cite{harmaClassificationTimefrequencyRegions2010,harmaClassificationTimefrequencyRegions2011,harmaStereoAudioClassification2011} collectively define a new paradigm of classifying stereo field whereby stereo mixtures are modelled as combinations of centre, panned, and uncorrelated sources; Although, they specifically state that most music signals will fall into a class that is very difficult to identify and is therefore omitted from the model.
\cite{prestiVisualizationManipulationStereophonic2014} uses the \gls{iid} and \gls{ipd} of an audio signal to derive bin-level measurements of \enquote{panning position} and \enquote{panning clarity}, which are subsequently used for sonogram visualization and manipulation.
Similar to the phase-based features defined in \cite{tardieuProductionEffectAudio2011}, there are a group of algorithms based on the emulation of a studio goniometer - \cite{ziemerAcousticFeaturesRecording2021,ziemerGoniometersArePowerful2023} use both a goniometer box-counting technique and inter-channel correlation analysis to derive a time-varying measurement of spaciousness. Interestingly though, \cite{stirnatSpaciousnessMusicTonmeisters2019} found that subjective impressions of spaciousness in various musical styles were largely contradictory to the results of goniometer analysis.

\gls{dirac} is a technique that converts ambisonic audio to a representation of directionality and diffuseness \cite{pulkkiSpatialSoundReproduction2007}. While it can be used on stereo audio, the conversion from stereo to ambisonic assumes the use of speakers rather than headphones \cite{laitinenConvertingTwochannelStereo2014}.

More recently a \gls{ml}-based technique for disentangling spatial information has been demonstrated with a use case of upmixing stereo to multichannel \cite{yangUpmixingStyleTransfer2022}; The input to the system is a combination of multichannel audio and it's associated stereo downmix, which is then re-upmixed by a neural network. A latent vector formed by running multichannel audio through one half of the network is demonstrated to represent the spatial information contained in the original signal. 



\section{State of practice}

Even though all the audio formats mentioned contain spatial information, the industry uses the term \enquote{spatial audio} to mean multichannel, ambisonic, or object-based. The main industrial use for rendered audio is compression and transmission/distribution, and the formats that have emerged generally allow for all 3 audio types to be used as input, even if they will not necessarily remain in the same format when a file is produced.

The most common industry format is Dolby Atmos - originally designed for cinemas and other such very large speaker arrays, there is now focus in rendering it in headphones for consumer users. Headphone rendering sometimes entails the use of head tracking to make it appear as though audio is always coming from a specific direction, although this is generally limited to specific hardware platforms where tracking data is available. Even though Atmos is being heavily pushed by big industry players, adoption by consumers is slow. An international standard spatial format exists in the form of MPEG-H \cite{mpeg-h}, (and its branded version \enquote{Sony 360}), but adoption of them is even lower. DTS:X is another proprietary spatial audio format, although its usage is generally limited to films released on physical media \cite{publishedDTSWhatIt2021}. Use of the industry formats is sometimes problematic due to patent licensing, and as a response the Alliance for Open Media has also specified a \enquote{royalty-free} standard for immersive audio named \gls{iamf} \cite{ImmersiveAudioModel}.

\section{Research Methods}
The papers mentioned in \cref{sec:papers} used (or will use) the following research methods:

\begin{itemize}
    \item \textbf{\nameref{paper:A}}: Several \gls{dsp} techniques for determining audio spatial features were performed on a small dataset that had predictable musical characteristics. The different techniques were compared using a new objective metric.
    \item \textbf{\nameref{paper:B}}: A perceptual test was used to determine listeners' reactions to carefully constructed audio stimuli. Standard statistical analysis was performed on the raw test results, and one particular \gls{dsp} technique (\gls{ph} \cite{gomezMusicSimilaritySystems2008} which was assessed in \nameref{paper:A}) was used to explain the perceptual results.
    \item \textbf{\nameref{paper:C}}: Following on the small selection of \gls{dsp} techniques assessed in \nameref{paper:A}, a \gls{slr} following the \gls{prisma} standard was performed to gather all currently available methods for determining spatial features of stereophonic musical audio.
    \item \textbf{\nameref{paper:D}}: Based on the results of the previous three papers, several types of \gls{sip} will be extracted from a very large dataset of recordings, and basic unsupervised machine-learning will be applied to find patterns within the \gls{sip} population.
\end{itemize}


\section{Preliminary results and conclusions}

\nameref{paper:A} revealed that although there are several \gls{dsp} techniques specifically designed for the extraction of recording and panning features, they are all limited in their own ways. One particular technique (\gls{spf} from \cite{tzanetakisStereoPanningFeatures2007}) was particularly promising as a candidate for use in \nameref{paper:D}, and the others could be used with some modifications. Overall it highlighted the importance of choosing both a panning feature \gls{dsp} technique and a distance metric that are appropriate to the overarching task.

\nameref{paper:B} concluded that when listeners are presented with voice overlayed on top of music, the best panning position for the voice was not always in the default centre position. The reasons for this were unclear due to a shortage of data, but the paper did demonstrate a real-world use of \gls{sip} extraction (using a algorithm trialed in \nameref{paper:A}). Most relevant to the research questions presented in \cref{sec:problem_statement}, it revealed that an objective measurement of panning can show detail that is not easily discernable by listening, such as very slight panning away from the centre.

\nameref{paper:C} revealed several more techniques that can be used to extract panning features. Although, given the high number of studies that were assessed for the review, there were relatively few very-relevant ones. The \gls{slr} also revealed that while \gls{ml} has been used on spatial features in the past, both \textit{unsupervised} \gls{ml} and very large music corpora remain fairly unexplored.

Overall the results so far have highlighted the known technique for extracting panning features from music, demonstrated the musical features that the major ones are able to reveal, and highlighted gaps in current research that \nameref{paper:D} aims to fill.

\section{Remaining work for licentiate}
Initially there is a need to find a new publication venue for \nameref{paper:B}, and this work is already underway. \nameref{paper:D} therefore comprises the bulk of the remaining work. Firstly it involves setting up a system that can gather \glspl{sip} en-mass from real recorded music. Secondly there will be an exploratory phase where many of the standard unsupervised learning techniques will be attempted to discover which are suitable for \gls{sip} data. The results of the \gls{ml} will then need to be assessed in some manner. One option is to correlate them with some known music metadata to reveal which other features are linked to panning. Another option is to perform listening test in order to assess is the \gls{ml} worked as intended.

\printglossaries

\printbibliography

\clearpage  % Forces BibTeX to process the entire document

\end{document}